% ============================================================
% CHAPTER 8 — External Validation
% ============================================================

\chapter{External Validation}
\label{chap:external_validation}

Chapter~\ref{chap:internal_validation} established that the proposed
framework's behaviour is attributable to the combination of \ac{PPO}
and \ac{DCR} shielding, on the Loan Application benchmark used
throughout its development. That chapter, together with
Chapter~\ref{chap:results}, answers whether the method is internally
correct: whether it does what it is designed to do, on the benchmarks
used to design it. This chapter asks a different question: whether
that correctness generalises to process structures that played no role
in shaping the method, that is, whether the findings of
Chapters~\ref{chap:results} and~\ref{chap:internal_validation} hold
under genuine external validity rather than only on the benchmarks
where the reward function, the shield, and the training
hyperparameters were calibrated.

Section~\ref{sec:external_benchmark_grid} introduces the benchmark
suite used for this purpose and the criteria for its selection.
Section~\ref{sec:external_design} describes the experimental design,
in particular the four conditions compared and the rationale for
allocating training budget unevenly across the suite.
Section~\ref{sec:external_results} reports the results: a replication
of the compliance finding of contribution~C1 on graphs never seen
during development, and a scalability frontier that neither the exact
solver nor the unmasked learning agent can cross, but that shielding
alone can. Section~\ref{sec:external_policy} reports a secondary
analysis of the sensitivity of training convergence to the
cost/duration assignment policy. Section~\ref{sec:external_discussion}
discusses the scope of these conclusions.

% ============================================================
\section{Benchmark Selection: Mined DCR Graphs}
\label{sec:external_benchmark_grid}

External validity requires a benchmark suite that satisfies two
conditions that none of the graphs used so far in this thesis satisfy
simultaneously. First, the process structure must be of independent
provenance, not authored or curated for this thesis. Second, no
parameter of the method, the reward shaping, the shield, the
entropy coefficient, the step penalty, may have been adjusted in
response to observed behaviour on that structure. The Loan Application
benchmark, while derived from a real event log, fails the second
condition: it was part of the same development cycle in which the
shield and the reward formulation were calibrated
(Chapter~\ref{chap:internal_validation}), and its cost values were
assigned by hand rather than by a fixed, reproducible procedure.

The benchmark suite used in this chapter is drawn from a set of \ac{DCR}
graphs mined by an independent thesis~\cite{TobiasDCRDiscover} from
real-world event logs via \ac{DCR}-discover, spanning processes such as
BPI Challenge incidents, loan applications, and bank transactions. The
discovered models are published directly under \texttt{datasets/models}
of that thesis's benchmarking repository,\footnote{Benchmark source and
discovered models: \url{https://github.com/s205358/dcr-js/tree/s205358/bench};
the exact modifications relative to the baseline DCR-js architecture can
be reviewed via \url{https://github.com/s205358/dcr-js/compare/s205358/refactor...s205358/bench}.}
which also provides a script for downloading the corresponding raw
event logs, with execution instructions in that repository's README.
None of these graphs were inspected before the present chapter, and the
cost/duration value assigned to each activity is produced by a fixed,
seeded random procedure (Section~\ref{sec:external_policy}),
removing manual judgement from the benchmark construction.

\paragraph{Selection grid.}
Six graphs were selected to span the size, density, and constraint-type
dimensions along which a declarative process model can vary. Size
follows the small (1--10 activities), medium (11--30), and large (30+)
categories associated with cognitive-overload thresholds in the
process-modelling literature. Density follows the formal definition of
Abbad-Andaloussi et al.~\cite{AbbadAndaloussi2023Complexity}: for a
graph $G$ with weakly-connected components $\mathrm{Comp}(G)$, each
with activity set $A_c$ and constraint set $C_c$,
\[
  \mathrm{Density}(G) = \max_{c \,\in\, \mathrm{Comp}(G)} \frac{|C_c|}{|A_c|},
\]
rather than a flat ratio of constraints to activities over the whole
graph, which would understate the density of any graph with
disconnected fragments. Table~\ref{tab:external_grid} reports the
resulting grid.

\begin{table}[H]
\centering
\caption{Benchmark suite for external validation: mined \ac{DCR}
graphs spanning the size/density/constraint-type grid. Density follows
Abbad-Andaloussi et al.~\cite{AbbadAndaloussi2023Complexity}. COP status
reports whether the exact CSP+COP solver of Diaz et
al.~\cite{Diaz2024Pareto-OptimalModels} converges within a 300\,s
timeout.}
\label{tab:external_grid}
\begin{tabular}{lrrlll}
\toprule
\textbf{Graph} & \textbf{Activities} & \textbf{Density} &
\textbf{Size} & \textbf{Dominant type} & \textbf{COP (300\,s)} \\
\midrule
04 (BPI13 Incidents)     & 4   & 1.75  & small  & mixed           & solved \\
01 (Artificial 0-noise)  & 8   & 1.60  & small  & exclusion-heavy & solved \\
05 (Synthetic Review Lg) & 14  & 6.00  & medium & condition-heavy & solved (2 points) \\
03 (BPI20 RFP)           & 19  & 10.32 & medium & mixed           & solved \\
11 (BPI19)               & 42  & 14.52 & large  & exclusion-heavy & solved (2 points) \\
09 (Large Bank Txn)      & 113 & 12.17 & large  & exclusion-heavy & \textbf{timeout} \\
\bottomrule
\end{tabular}
\end{table}

A structural Breadth-First-Search analysis of the \ac{DCR} marking
semantics, mirroring exactly the acceptance condition of the CSP+COP
model, shows that all six graphs accept the empty trace
($\mathit{Pending} \cap \mathit{Included} = \emptyset$ at the initial
marking): they are mined from event logs whose discovered pending-set
structure is trivially satisfied. This matches a risk already
identified in earlier benchmark selection
(Section~\ref{sec:graph_selection_criterion}): graphs mined directly
from event logs frequently lack the non-trivial response structure
that makes multi-objective optimisation meaningful. Forcing at least
one real transition, exactly as the CSP+COP model does by construction
since it cannot represent the empty trace, four of the six graphs (04,
01, 03, 11) admit a trivial one-step accepting continuation: any
legal first action already satisfies acceptance. These four are
retained as \emph{controls} rather than as headline results: they
confirm that the framework behaves correctly when the optimisation
problem is degenerate, but do not exercise the shielding mechanism
under a genuine multi-step trade-off. Graph 05 is the only one of the
six whose non-trivial minimum accepting trace requires a multi-step
chain (six steps) while still admitting an exact CSP+COP solution, and
is therefore treated as the anchor result of this chapter. Graph 09's
non-trivial minimum accepting trace requires 51 sequential steps,
found by the structural BFS in 0.31\,s, while the CSP+COP solver does
not converge within 300\,s; it is the scale case of this chapter. No
partial or incumbent solution is available for comparison on this
graph either: the solver's enumeration procedure (which iteratively
finds an improving feasible trace, then branches to require a further
improvement) was re-run with the same 300\,s budget while instrumented
to retain whatever candidate traces had been found if the deadline was
reached mid-enumeration, rather than discarding them on a hard kill as
the main sweep does. Even the very first feasible trace, prior to any
optimisation step, did not return within several minutes beyond the
budget; the solver does not reach the search depth required even to
begin enumerating candidates on this graph, so there is no incumbent
front to report, exact or partial.

This outcome is not unexpected given the scale of the graph relative to
where the exact method itself is known to degrade: Diaz et
al.~\cite{Diaz2024Pareto-OptimalModels} report that their CSP+COP
approach begins to fail in practice at approximately 45 activities,
and graph 09's non-trivial minimum accepting trace, at 51 sequential
steps, already exceeds that range before any optimisation is
attempted. As a purely structural reference point, not a claim of
optimality, this 51-step trace (found by the BFS reachability search,
independently of MiniZinc) costs 2877 and has a duration of 3212 under
the cost/duration assignment of
Section~\ref{sec:external_policy}; it is reported here only
to give a concrete sense of scale to the otherwise empty CSP+COP
column for this graph, and is excluded from any Pareto or hypervolume
comparison in the remainder of this chapter.

% ============================================================
\section{Experimental Design}
\label{sec:external_design}

Four conditions are compared on each graph, extending the ablation of
Section~\ref{sec:shield_contribution} from a single internal benchmark
to the full grid.

\begin{description}
    \item[COP.] The exact CSP+COP solver, reported where it converges
    within the 300\,s budget (Table~\ref{tab:external_grid}). Serves as
    ground truth.
    \item[\textnormal{\texttt{saferl}}.] \ac{PPO} with \ac{DCR}
    shielding, the full framework, identical configuration to
    Chapters~\ref{chap:results} and~\ref{chap:internal_validation}.
    \item[\textnormal{\texttt{shield\_only}}.] The non-learning random
    valid-action baseline of Section~\ref{sec:shield_contribution},
    isolating how much of the result is attributable to structural
    restriction alone.
    \item[\textnormal{\texttt{rl\_only}}.] \ac{PPO} with the shield
    disabled, the same ablation arm as in
    Section~\ref{sec:shield_contribution}.
\end{description}

\paragraph{Budget allocation.}
Training budget was allocated unevenly across the grid, in proportion
to how much signal each graph can provide. The four control graphs
(04, 01, 03, 11) were trained for two representative
$(\alpha,\beta)$ pairs at 25\,000 steps; local pilot runs showed the
full compliance pattern emerges within a few thousand steps on these
small action spaces, so the full six-pair, 100\,000-step protocol used
elsewhere in this thesis would add wall-clock cost without additional
signal. Graph 05, the anchor result, received the full protocol: six
weight pairs, 100\,000 steps per pair, for \texttt{saferl} and
\texttt{rl\_only}, plus a \texttt{shield\_only} run capped at the same
100\,000-step budget. Graph 09 was trained only under
\texttt{shield\_only}: a preliminary diagnostic run of \texttt{saferl}
at 50\,000 steps showed a 0\,\% acceptance rate with the illegal-action
ratio plateaued at 98--100\,\% with no improving trend, attributable to
the size of the action space rather than to the length of the
required trace (Section~\ref{sec:external_discussion}); repeating this
at the full budget would not be expected to change the qualitative
outcome and was not pursued. \texttt{rl\_only} was likewise not run on
graph 09, for the same reason.

\paragraph{Consistency with the COP ground truth.}
Cost and duration values are assigned per activity by a fixed,
seeded procedure (uniform sampling within the same ranges used for the
internal benchmarks), so that the \ac{RL} environment and the CSP+COP
model optimise over identical numerical values, an apples-to-apples
comparison that the synthetic and hand-assigned costs of the internal
benchmarks do not by themselves guarantee.

\paragraph{Empty-trace consistency.}
Because all six graphs structurally accept the empty trace
(Section~\ref{sec:external_benchmark_grid}), and the CSP+COP model
cannot represent it, the \ac{RL} environment must enforce the same
restriction for the comparison to be valid. This is already the case
by construction: episode termination is evaluated only inside the
environment's step function, never at reset, so the agent cannot
receive an accepting signal without having executed at least one
action. No change to the environment was required to establish this
consistency.

% ============================================================
\section{Results}
\label{sec:external_results}

\subsection{Scalability Frontier}
\label{sec:external_scalability}

Figure~\ref{fig:external_accept_vs_size} reports the accept rate in the
final 20\,\% of training (or evaluation, for \texttt{shield\_only}) for
each condition, as a function of graph size.

\begin{figure}[H]
    \centering
    \includegraphics[width=0.95\linewidth]{fig2_accept_rate_vs_size.png}
    \caption{Accept rate (final 20\,\% of episodes) versus graph size,
    by condition, across the six-graph benchmark grid. Markers,
    linestyles, and a small horizontal jitter distinguish the three
    conditions even where they coincide at 1.0. \texttt{rl\_only} was
    not run on graph 09 (marked explicitly), since the diagnostic run
    already indicated it would fail at least as severely as
    \texttt{saferl}.}
    \label{fig:external_accept_vs_size}
\end{figure}

\texttt{shield\_only} sustains 100\,\% acceptance across the entire
grid, including the 113-activity graph where the exact solver times
out. \texttt{saferl} also sustains 100\,\% acceptance up to 42
activities, but collapses to 0\,\% at 113 activities, approximately the
same scale at which the CSP+COP solver stops converging. The
non-learning, structurally restricted baseline is therefore the only
condition in this study that scales to where exact enumeration becomes
infeasible; the learning component adds value only within the range
where it can still compete with the size of the action space (up to 42
activities in this study), not beyond it.

The six-graph grid, however, has no test point between 42 and 113
activities, so it can only establish that both methods fail somewhere
within that 71-activity gap, not that they fail at a shared threshold.

\subsection{Controlled Scaling Diagnostic}
\label{sec:external_scaling_diagnostic}

To localise the crossover more precisely, both methods were additionally
run on a series of controlled, synthetically generated chain-plus-noise
\ac{DCR} graphs at five intermediate sizes (20, 40, 60, 80, and 100
events), generated with the same seed and generator for both methods so
that each pair of runs sees an identical graph. Unlike the mined-graph
grid of Table~\ref{tab:external_grid}, these graphs are not drawn from
real event logs and therefore carry no claim to external validity; their
purpose is resolution, not generalisation, complementing the mined grid
rather than replacing it. The \ac{RL} side was trained at a reduced,
diagnostic step budget (20,000 steps, baseline weights $\alpha=\beta=0$)
rather than the full per-configuration budget used elsewhere in this
thesis, since the goal is to locate the wall, not to characterise
optimisation quality near it.

\begin{table}[H]
\centering
\caption{Controlled scaling diagnostic: CSP+COP convergence and
\ac{PPO}-with-shield accept rate (final 20\,\% of episodes) on synthetic
chain-plus-noise graphs of increasing size. CSP+COP times out at the
same 300\,s budget used in Table~\ref{tab:external_grid}.}
\label{tab:scaling_diagnostic}
\begin{tabular}{rrrl}
\toprule
\textbf{Events} & \textbf{CSP+COP} & \textbf{PPO accept rate} & \textbf{CSP+COP solve time} \\
\midrule
20  & solved  & 100\,\%  & 0.005\,s \\
40  & solved  & 100\,\%  & 0.015\,s \\
60  & timeout & 7.7\,\%  & --- \\
80  & timeout & 0\,\%    & --- \\
100 & timeout & 0\,\%    & --- \\
\bottomrule
\end{tabular}
\end{table}

Both methods succeed at 40 events and both fail by 60, narrowing the
shared crossover from the 71-activity gap of the mined grid to a
20-activity window. The fact that this pattern holds on a graph family
unrelated to the one mined-graph grid, generated independently and
evaluated under a reduced training budget, is itself evidence that the
shared failure point is not an artefact of the specific six mined
graphs: the exact solver's constraint-explosion wall and the unmasked
policy's action-space wall sit in the same narrow region regardless of
which graph family is used to probe it. This finer-grained result should
be read as a diagnostic complement to Section~\ref{sec:external_scalability},
not as a substitute for it: it narrows \emph{where} the shared wall is,
while the mined-graph grid establishes that the underlying findings of
this chapter hold on graphs with genuine external provenance.

\subsection{Replication of the Compliance Finding (C1)}
\label{sec:external_compliance_replication}

Figure~\ref{fig:external_illegal_deciles} reports the illegal-action
rate per training decile for \texttt{saferl} and \texttt{rl\_only},
pooled across weight configurations, on each of the six graphs.

\begin{figure}[H]
    \centering
    \includegraphics[width=\linewidth]{fig1_illegal_rate_deciles.png}
    \caption{Illegal-action rate per training decile, \texttt{saferl}
    versus \texttt{rl\_only}, across the six-graph benchmark grid.}
    \label{fig:external_illegal_deciles}
\end{figure}

On the four control graphs, the gap between conditions is small: with
4--19 activities and a trivially satisfiable acceptance condition,
random action choice already has a reasonable probability of being
legal, so the shield has comparatively little work to do. On graph 05,
the only control-free graph with a genuine multi-step trade-off and a
14-activity action space, the gap is large and closely replicates the
pattern observed internally on the Loan Application benchmark
(Table~\ref{tab:external_replication}).

\begin{table}[H]
\centering
\caption{Replication of the compliance finding of
Section~\ref{sec:rl_vs_safe_rl} on an external benchmark (graph 05,
final training decile).}
\label{tab:external_replication}
\begin{tabular}{lcc}
\toprule
\textbf{Metric} & \textbf{Loan Application (internal)} & \textbf{Graph 05 (external)} \\
\midrule
\texttt{saferl} illegal rate  & 1.0\,\%  & 1.77\,\% \\
\texttt{rl\_only} illegal rate & 35.8\,\% & 50.7\,\% \\
\bottomrule
\end{tabular}
\end{table}

The replication of this gap on a graph never inspected during the
development of the shield or the reward formulation is the central
external-validity evidence for contribution~C1: a flat penalty signal
is not sufficient to obtain compliant behaviour, and the gap is not an
artefact specific to the Loan Application benchmark.

\subsection{Pareto Recovery Against the Exact Solver}
\label{sec:external_pareto}

Figure~\ref{fig:external_pareto} compares, for the five graphs where
the CSP+COP solver converges, the non-dominated subset of
\texttt{saferl}'s accepting episodes in the final 20\,\% of training
against the exact Pareto front.

\begin{figure}[H]
    \centering
    \includegraphics[width=\linewidth]{fig3_pareto_vs_cop.png}
    \caption{Non-dominated front recovered by \texttt{saferl} (green,
    connected) against the exact CSP+COP front (hollow stars, drawn
    behind so a coincident point remains visible), for the five graphs
    with a converged exact solution. Dominated accepting episodes are
    shown in light grey for reference.}
    \label{fig:external_pareto}
\end{figure}

In all five graphs, the non-dominated subset of \texttt{saferl}'s
accepting episodes coincides exactly with the CSP+COP front. This is
verified explicitly rather than read off the figure alone: on graph
05, where the exact front has two points, $\{(335,410),(339,387)\}$,
the non-dominated subset of \texttt{saferl}'s accepting episodes is
exactly this set, no more and no fewer points. On the remaining four
graphs, where the exact front is a single point, the same holds.

Table~\ref{tab:external_front_quality} reports this agreement using
four standard multi-objective quality indicators, rather than only the
hypervolume already given in Table~\ref{tab:external_master}: the
Averaged Hausdorff Distance (AHD) between the recovered and exact
fronts, the Purity of the recovered front (the fraction of its points
that remain non-dominated once pooled with the exact front), and the
Delta spacing metric (how evenly the recovered points populate the
front). Objectives are min--max normalised to $[0,1]$ per graph before
computing distances, so cost and duration contribute on comparable
scales despite their very different absolute magnitudes across graphs.

\begin{table}[H]
\centering
\caption{Front-quality indicators of \texttt{saferl}'s recovered front
against the exact CSP+COP front, for the five graphs where the solver
converges. AHD$=0$ and Purity$=1$ throughout confirm the exact-match
claim quantitatively. Delta is undefined ("--") for single-point
fronts.}
\label{tab:external_front_quality}
\begin{tabular}{lccccc}
\toprule
\textbf{Graph} & \textbf{Points} & \textbf{AHD} & \textbf{Purity} & \textbf{Delta} \\
\midrule
04 & 1 & 0.0 & 1.0 & -- \\
01 & 1 & 0.0 & 1.0 & -- \\
05 & 2 & 0.0 & 1.0 & 0.0 \\
03 & 1 & 0.0 & 1.0 & -- \\
11 & 2 & 0.0 & 1.0 & 0.0 \\
\bottomrule
\end{tabular}
\end{table}

\paragraph{An apparent anomaly in \texttt{rl\_only}.}
The hypervolume of \texttt{rl\_only}'s accepted episodes, computed
against the same reference point as the exact front, is larger than
the CSP+COP hypervolume on graphs 05 and 11 (239{,}070 versus 34{,}655,
and 1{,}392 versus 348, respectively), which would suggest, taken at
face value, that the unshielded agent finds a better front than the
exact solver, an impossibility if the solver's front is truly optimal.
Inspection of the underlying episodes resolves the apparent
contradiction: these are single- or few-step episodes reaching
artificially low cost and duration through illegal transitions that
corrupt the marking, the same phenomenon, here on a different graph,
identified as the \texttt{Event\_10} exploit in
Section~\ref{sec:shield_contribution}. They are not better solutions;
they are not solutions to the modelled process at all. This
reinforces, on an independent benchmark, that hypervolume alone is an
insufficient quality measure for an unshielded agent, and that
executed-trace validity must be reported alongside it.

\subsection{Summary}
\label{sec:external_summary}

Table~\ref{tab:external_master} consolidates the results of this
section.

\begin{table}[H]
\centering
\caption{Summary of external-validation results across the six-graph
benchmark grid. Illegal rate is reported for the final training decile.
Hypervolume is computed against a fixed reference point per graph;
\texttt{rl\_only}'s hypervolume on graphs 05 and 11 reflects invalid
episodes, not superior solutions (Section~\ref{sec:external_pareto}).}
\label{tab:external_master}
\begin{tabular}{lccccc}
\toprule
\textbf{Graph} & \textbf{\texttt{saferl} accept} & \textbf{\texttt{rl\_only} accept}
& \textbf{\texttt{shield\_only} accept} & \textbf{\texttt{saferl} illegal\,\%}
& \textbf{\texttt{rl\_only} illegal\,\%} \\
\midrule
04 & 100\,\%   & 100\,\% & 100\,\% & 0.18 & 0.13 \\
01 & 100\,\%   & 100\,\% & 100\,\% & 0.12 & 0.30 \\
05 & 100\,\%   & 100\,\% & 100\,\% & 1.77 & 50.7 \\
03 & 100\,\%   & 100\,\% & 100\,\% & 3.76 & 2.69 \\
11 & 99.98\,\% & 100\,\% & 100\,\% & 1.68 & 5.49 \\
09 & \textbf{0\,\%} & not run & \textbf{100\,\%} & 98.4 & -- \\
\bottomrule
\end{tabular}
\end{table}

% ============================================================
\section{Cost/Duration Assignment Policy Sensitivity}
\label{sec:external_policy}

A secondary question, raised independently during the design of this
benchmark suite, is whether the procedure used to assign synthetic
cost and duration values to each activity, here a fixed, seeded
random draw, affects training convergence or the geometry of the
recovered front. Two structured alternatives were compared against the
random assignment on graph 05, holding the duration draw identical
across all three (only the cost value differs, as a function of
duration):

\begin{description}
    \item[Linear.] Cost increases with duration (a process in which
    spending more time also costs more, e.g.\ billed by the hour).
    \item[Inverse.] Cost decreases with duration (a process in which
    paying more buys a faster execution, e.g.\ an expedited service).
\end{description}

\begin{table}[H]
\centering
\caption{Effect of the cost/duration assignment policy on training
convergence and recovered front, graph 05, \texttt{saferl}. All three
policies share an identical duration draw; only the cost mapping
differs. Accept rate and illegal rate are measured over the final
20\,\% / final decile of training, pooled across the six weight pairs.}
\label{tab:external_policy}
\begin{tabular}{lccccc}
\toprule
\textbf{Policy} & \textbf{Accept rate} & \textbf{Illegal rate (final decile)} & \textbf{Pareto points} & \textbf{Cost range} & \textbf{Duration range} \\
\midrule
Random  & 100\,\% & 3.53\,\% & 1 & 371--371 & 350--350 \\
Linear  & 100\,\% & 4.83\,\% & 1 & 295--295 & 350--350 \\
Inverse & 100\,\% & 2.78\,\% & 2 & 327--335 & 350--360 \\
\bottomrule
\end{tabular}
\end{table}

\begin{figure}[H]
  \centering
  \includegraphics[width=0.75\textwidth]{Pictures/policy_convergence.png}
  \caption{Illegal-action rate per training decile under the three
  cost/duration assignment policies, graph 05, \texttt{saferl}, pooled
  across the six weight pairs. The three curves are visually
  indistinguishable throughout training.}
  \label{fig:policy_convergence}
\end{figure}

\begin{figure}[H]
  \centering
  \includegraphics[width=0.75\textwidth]{Pictures/policy_front_geometry.png}
  \caption{Recovered front and dominated episode cloud under the three
  cost/duration assignment policies, graph 05, \texttt{saferl}, final
  20\,\% of training. The dominated cloud reflects the cost--duration
  correlation imposed by each policy (ascending diagonal under Linear,
  descending under Inverse, no pattern under Random); the recovered
  front itself lands in a similar region under all three.}
  \label{fig:policy_front_geometry}
\end{figure}

Convergence speed and final compliance are stable across the three
policies: the illegal-rate curves in
Figure~\ref{fig:policy_convergence} are visually indistinguishable
throughout training, and the final illegal rate stays within a narrow
band (2.78--4.83\,\%) regardless of assignment policy. The recovered
front also lands in a comparable cost--duration region under all three
(Figure~\ref{fig:policy_front_geometry}), though it is a single point
under Random and Linear and two points under Inverse --- a difference
of degree, not of qualitative behaviour, given the narrow ranges
involved. What does change visibly is the geometry of the
\emph{dominated} episode cloud, which traces an ascending diagonal under
Linear, a descending diagonal under Inverse, and no discernible pattern
under Random: this is the direct, expected consequence of each policy's
per-activity cost--duration correlation, not evidence that one policy
makes the optimisation problem harder. Taken together, these results
support the conclusion that the choice of cost/duration assignment
procedure is a property of the benchmark's difficulty surface, not of
the learning dynamics: it shapes which trade-offs exist for the agent to
find, but not how quickly or reliably the agent finds them.

% ============================================================
\section{Discussion and Scope}
\label{sec:external_discussion}

This chapter establishes that the two central findings of
Chapters~\ref{chap:results} and~\ref{chap:internal_validation}, exact
Pareto recovery and the necessity of shielding for compliant behaviour,
generalise to process structures of independent provenance, never
inspected during the development of the method. It additionally
exposes a scalability frontier that the internal benchmarks, all under
20 events, could not reveal: the learning component's contribution is
bounded by the size of the action space it must search without
guidance, and beyond that bound the simplest possible compliant
baseline, a shield with no learning at all, is the only condition in
this study that continues to produce valid executions where the exact
solver and the unmasked learning agent both fail.

This finding refines, rather than weakens, the practical claim of
shielding. The version supported by this chapter's evidence is not
that the shielded \ac{RL} framework outscales exact constraint solving,
the results of Section~\ref{sec:external_scalability} do not support
that claim beyond 42 activities in this study, but that shielding by
itself, independent of whether it is paired with a learning component,
is sufficient to keep producing structurally valid traces in exactly
the regime where exact enumeration becomes computationally infeasible.
Recovering the learning component's advantage at that scale, that is,
not merely valid but optimised traces beyond a few dozen activities,
would require restricting the policy's action sampling to the
currently enabled events directly (action masking), rather than relying
on the shield to intercept illegal samples after the fact; this is left
as future work and is not evaluated here, since it would constitute a
change to the learning method itself rather than to its validation.

Two further scope limitations apply. First, the comparison against
Kirchdorfer et al.'s evaluation design
(Section~\ref{sec:loanapp_capacity}) is not part of this chapter's
external-validity claim: it is run on the same Loan Application
benchmark used internally, and is best understood as a methodological
replication, not as evidence of generalisation to unseen structure. An
exhaustive search of the literature (Section~\ref{sec:related_work})
found no published benchmark combining reinforcement learning with
\ac{DCR} graphs against which the role-conditioned, capacity-constrained
extension of this thesis could be externally compared; the closest
related work, Agarwal et al.~\cite{Agarwal2022GoalOriented}, addresses
a related but distinct problem (goal-oriented activity recommendation
under a directly-follows process model, not a declarative one) and
independently confirms, via its use of action masking, that the same
action-space bottleneck identified in
Section~\ref{sec:external_scalability} is a recognised issue in this
family of problems generally, not one specific to this thesis's
implementation. Second, the four control graphs (04, 01, 03, 11) are
retained in this chapter for completeness and transparency, since
discarding them silently would understate how much of the mined-graph
benchmark population is structurally trivial, but they do not, by
construction, exercise the framework under genuine multi-step
optimisation, and the headline claims of this chapter rest on graphs
05 and 09 alone.
